\section{Vue.js在智慧水利平台中的实践}\label{vue.jsux5728ux667aux6167ux6c34ux5229ux5e73ux53f0ux4e2dux7684ux5b9eux8df5}

本节将探讨Vue.js在智慧水利平台开发中的实际应用，包括常见功能模块的实现方案、最佳实践和性能优化策略。

\subsection{6.1
智慧水利平台功能模块}\label{ux667aux6167ux6c34ux5229ux5e73ux53f0ux529fux80fdux6a21ux5757}

智慧水利平台通常包含多个功能模块，每个模块都可以利用Vue.js的特性进行高效开发。

\subsubsection{6.1.1
GIS地图与水文监测}\label{gisux5730ux56feux4e0eux6c34ux6587ux76d1ux6d4b}

GIS地图是智慧水利平台的核心组件，用于展示水文站点分布和监测数据。

\begin{lstlisting}
<!-- WaterMonitoringMap.vue -->
<template>
  <div class="map-container">
    <div id="gisMap" class="map-view"></div>
    
    <div class="map-controls">
      <div class="layer-control">
        <h3>图层控制</h3>
        <el-checkbox-group v-model="activeLayers">
          <el-checkbox label="stations">水文站点</el-checkbox>
          <el-checkbox label="reservoirs">水库</el-checkbox>
          <el-checkbox label="rivers">河流</el-checkbox>
          <el-checkbox label="rainfall">降雨量</el-checkbox>
          <el-checkbox label="watershed">流域范围</el-checkbox>
        </el-checkbox-group>
      </div>
      
      <div class="time-control">
        <h3>时间控制</h3>
        <el-date-picker 
          v-model="currentTime" 
          type="datetime" 
          placeholder="选择日期时间"
          format="yyyy-MM-dd HH:mm"
          @change="updateTimeData">
        </el-date-picker>
      </div>
    </div>
    
    <station-detail-panel 
      v-if="selectedStation" 
      :station="selectedStation"
      :time="currentTime"
      @close="closeStationDetail">
    </station-detail-panel>
  </div>
</template>

<script>
import L from 'leaflet'
import 'leaflet/dist/leaflet.css'
import StationDetailPanel from './StationDetailPanel.vue'
import { mapState, mapActions } from 'vuex'

export default {
  components: {
    StationDetailPanel
  },
  data() {
    return {
      map: null,
      activeLayers: ['stations', 'rivers'],
      currentTime: new Date(),
      layerGroups: {},
      selectedStation: null
    }
  },
  computed: {
    ...mapState('waterData', [
      'stations',
      'reservoirs',
      'watersheds',
      'rivers',
      'rainfallData'
    ])
  },
  watch: {
    activeLayers: {
      handler(newLayers) {
        this.updateLayerVisibility(newLayers);
      },
      deep: true
    },
    stations: {
      handler(newStations) {
        this.updateStationsLayer(newStations);
      },
      deep: true
    }
  },
  mounted() {
    this.initMap();
    this.fetchWaterData(this.currentTime);
  },
  methods: {
    ...mapActions('waterData', [
      'fetchWaterData', 
      'fetchStationDetail'
    ]),
    
    initMap() {
      this.map = L.map('gisMap').setView([30.5928, 114.3055], 8);
      
      // 添加底图
      L.tileLayer('https://{s}.tile.openstreetmap.org/{z}/{x}/{y}.png', {
        attribution: '&copy; <a href="https://www.openstreetmap.org/copyright">OpenStreetMap</a> contributors'
      }).addTo(this.map);
      
      // 初始化图层组
      this.layerGroups = {
        stations: L.layerGroup().addTo(this.map),
        reservoirs: L.layerGroup().addTo(this.map),
        rivers: L.layerGroup().addTo(this.map),
        rainfall: L.layerGroup(),
        watershed: L.layerGroup()
      };
      
      // 绑定地图事件
      this.map.on('zoomend', this.updateMapDisplay);
      this.map.on('moveend', this.updateMapDisplay);
    },
    
    updateLayerVisibility(activeLayers) {
      Object.keys(this.layerGroups).forEach(layerName => {
        const layerGroup = this.layerGroups[layerName];
        if (activeLayers.includes(layerName)) {
          if (!this.map.hasLayer(layerGroup)) {
            this.map.addLayer(layerGroup);
          }
        } else {
          if (this.map.hasLayer(layerGroup)) {
            this.map.removeLayer(layerGroup);
          }
        }
      });
    },
    
    updateStationsLayer(stations) {
      const stationsLayer = this.layerGroups.stations;
      stationsLayer.clearLayers();
      
      stations.forEach(station => {
        const marker = this.createStationMarker(station);
        marker.on('click', () => {
          this.handleStationClick(station.id);
        });
        stationsLayer.addLayer(marker);
      });
    },
    
    createStationMarker(station) {
      const statusColors = {
        normal: '#52c41a',
        warning: '#faad14',
        danger: '#ff4d4f'
      };
      
      const icon = L.divIcon({
        className: 'station-marker',
        html: `<div style="background-color: ${statusColors[station.status]}"></div>`,
        iconSize: [12, 12]
      });
      
      return L.marker([station.latitude, station.longitude], { icon });
    },
    
    async handleStationClick(stationId) {
      try {
        await this.fetchStationDetail(stationId);
        this.selectedStation = this.stations.find(s => s.id === stationId);
      } catch (error) {
        this.$message.error('获取站点详情失败');
      }
    },
    
    closeStationDetail() {
      this.selectedStation = null;
    },
    
    updateTimeData(time) {
      this.fetchWaterData(time);
    },
    
    updateMapDisplay() {
      // 根据当前缩放级别和可视区域更新数据
      const bounds = this.map.getBounds();
      const zoom = this.map.getZoom();
      
      // 可以在这里根据缩放级别和可视区域调整显示的数据粒度
      console.log('地图视图更新:', bounds, zoom);
    }
  },
  beforeDestroy() {
    if (this.map) {
      this.map.remove();
    }
  }
}
</script>

<style scoped>
.map-container {
  position: relative;
  width: 100%;
  height: 700px;
}

.map-view {
  width: 100%;
  height: 100%;
}

.map-controls {
  position: absolute;
  top: 10px;
  right: 10px;
  background-color: white;
  padding: 10px;
  border-radius: 4px;
  box-shadow: 0 2px 6px rgba(0, 0, 0, 0.2);
  z-index: 1000;
  min-width: 200px;
}

.layer-control, .time-control {
  margin-bottom: 15px;
}

.layer-control h3, .time-control h3 {
  margin-top: 0;
  margin-bottom: 8px;
  font-size: 16px;
}
</style>
\end{lstlisting}

\subsubsection{6.1.2
水库调度与模拟}\label{ux6c34ux5e93ux8c03ux5ea6ux4e0eux6a21ux62df}

水库调度模块用于水库水情监测和调度模拟，需要处理复杂的表单和实时数据。

\begin{lstlisting}
<!-- ReservoirOperation.vue -->
<template>
  <div class="reservoir-operation">
    <el-row :gutter="20">
      <el-col :span="8">
        <reservoir-info-card 
          :reservoir="currentReservoir"
          :real-time-data="realtimeData">
        </reservoir-info-card>
        
        <el-card class="operation-card">
          <div slot="header">
            <span>调度操作</span>
          </div>
          
          <el-form ref="operationForm" :model="operationForm" :rules="rules" label-width="100px">
            <el-form-item label="目标水位" prop="targetLevel">
              <el-input-number 
                v-model="operationForm.targetLevel" 
                :min="currentReservoir.deadWaterLevel" 
                :max="currentReservoir.normalWaterLevel"
                :step="0.1">
              </el-input-number>
              <span class="unit">米</span>
            </el-form-item>
            
            <el-form-item label="下泄流量" prop="outflowRate">
              <el-input-number 
                v-model="operationForm.outflowRate" 
                :min="0" 
                :max="5000"
                :step="10">
              </el-input-number>
              <span class="unit">立方米/秒</span>
            </el-form-item>
            
            <el-form-item label="调度时长" prop="duration">
              <el-input-number 
                v-model="operationForm.duration" 
                :min="1" 
                :max="72"
                :step="1">
              </el-input-number>
              <span class="unit">小时</span>
            </el-form-item>
            
            <el-form-item>
              <el-button type="primary" @click="submitOperation">提交调度方案</el-button>
              <el-button @click="runSimulation">运行模拟</el-button>
            </el-form-item>
          </el-form>
        </el-card>
      </el-col>
      
      <el-col :span="16">
        <el-card class="simulation-card">
          <div slot="header">
            <span>调度模拟结果</span>
          </div>
          
          <div v-if="simulationLoading" class="simulation-loading">
            <el-spinner></el-spinner>
            <p>模拟计算中...</p>
          </div>
          
          <div v-else-if="simulationResult" class="simulation-result">
            <el-tabs v-model="activeTab">
              <el-tab-pane label="水位变化" name="waterLevel">
                <water-level-chart 
                  :data="simulationResult.waterLevelData"
                  :warning-level="currentReservoir.floodWaterLevel">
                </water-level-chart>
              </el-tab-pane>
              
              <el-tab-pane label="入出库流量" name="flowRate">
                <flow-rate-chart :data="simulationResult.flowRateData"></flow-rate-chart>
              </el-tab-pane>
              
              <el-tab-pane label="蓄水量变化" name="storage">
                <storage-chart :data="simulationResult.storageData"></storage-chart>
              </el-tab-pane>
              
              <el-tab-pane label="下游影响" name="downstream">
                <downstream-impact-chart :data="simulationResult.downstreamData"></downstream-impact-chart>
              </el-tab-pane>
            </el-tabs>
            
            <div class="simulation-summary">
              <h3>模拟结果摘要</h3>
              <el-row :gutter="20">
                <el-col :span="8">
                  <div class="summary-item">
                    <div class="summary-label">最终水位</div>
                    <div class="summary-value">{{ simulationResult.finalWaterLevel.toFixed(2) }}米</div>
                  </div>
                </el-col>
                <el-col :span="8">
                  <div class="summary-item">
                    <div class="summary-label">累计下泄水量</div>
                    <div class="summary-value">{{ (simulationResult.totalOutflow / 10000).toFixed(2) }}万立方米</div>
                  </div>
                </el-col>
                <el-col :span="8">
                  <div class="summary-item" :class="{'warning': simulationResult.hasWarning}">
                    <div class="summary-label">预警信息</div>
                    <div class="summary-value">
                      {{ simulationResult.hasWarning ? simulationResult.warningMessage : '无预警' }}
                    </div>
                  </div>
                </el-col>
              </el-row>
            </div>
          </div>
          
          <div v-else class="no-simulation">
            <p>请设置调度参数并运行模拟</p>
          </div>
        </el-card>
      </el-col>
    </el-row>
  </div>
</template>

<script>
import ReservoirInfoCard from '@/components/reservoir/ReservoirInfoCard.vue'
import WaterLevelChart from '@/components/charts/WaterLevelChart.vue'
import FlowRateChart from '@/components/charts/FlowRateChart.vue'
import StorageChart from '@/components/charts/StorageChart.vue'
import DownstreamImpactChart from '@/components/charts/DownstreamImpactChart.vue'
import { mapState, mapActions } from 'vuex'

export default {
  components: {
    ReservoirInfoCard,
    WaterLevelChart,
    FlowRateChart,
    StorageChart,
    DownstreamImpactChart
  },
  props: {
    reservoirId: {
      type: String,
      required: true
    }
  },
  data() {
    return {
      operationForm: {
        targetLevel: 0,
        outflowRate: 0,
        duration: 24
      },
      rules: {
        targetLevel: [
          { required: true, message: '请输入目标水位', trigger: 'blur' }
        ],
        outflowRate: [
          { required: true, message: '请输入下泄流量', trigger: 'blur' }
        ],
        duration: [
          { required: true, message: '请输入调度时长', trigger: 'blur' }
        ]
      },
      simulationLoading: false,
      simulationResult: null,
      activeTab: 'waterLevel'
    }
  },
  computed: {
    ...mapState('reservoirs', [
      'reservoirList',
      'realtimeData'
    ]),
    
    currentReservoir() {
      return this.reservoirList.find(r => r.id === this.reservoirId) || {};
    }
  },
  watch: {
    currentReservoir: {
      handler(newReservoir) {
        if (newReservoir && newReservoir.id) {
          this.operationForm.targetLevel = newReservoir.currentWaterLevel || newReservoir.normalWaterLevel;
          this.operationForm.outflowRate = newReservoir.currentOutflow || 0;
        }
      },
      immediate: true
    }
  },
  created() {
    this.fetchReservoirData();
    this.startRealtimeDataPolling();
  },
  methods: {
    ...mapActions('reservoirs', [
      'fetchReservoirData',
      'fetchReservoirRealtimeData',
      'submitReservoirOperation',
      'simulateReservoirOperation'
    ]),
    
    startRealtimeDataPolling() {
      this.fetchReservoirRealtimeData(this.reservoirId);
      this.pollingTimer = setInterval(() => {
        this.fetchReservoirRealtimeData(this.reservoirId);
      }, 30000); // 每30秒更新一次实时数据
    },
    
    async submitOperation() {
      try {
        await this.$refs.operationForm.validate();
        
        await this.submitReservoirOperation({
          reservoirId: this.reservoirId,
          operation: { ...this.operationForm }
        });
        
        this.$message.success('调度方案提交成功');
        this.runSimulation();
      } catch (error) {
        console.error('提交调度方案失败:', error);
      }
    },
    
    async runSimulation() {
      try {
        await this.$refs.operationForm.validate();
        
        this.simulationLoading = true;
        this.simulationResult = await this.simulateReservoirOperation({
          reservoirId: this.reservoirId,
          parameters: { ...this.operationForm }
        });
        
        this.$message.success('模拟计算完成');
      } catch (error) {
        this.$message.error('模拟计算失败: ' + error.message);
        console.error('模拟计算失败:', error);
      } finally {
        this.simulationLoading = false;
      }
    }
  },
  beforeDestroy() {
    if (this.pollingTimer) {
      clearInterval(this.pollingTimer);
    }
  }
}
</script>

<style scoped>
.reservoir-operation {
  padding: 20px;
}

.operation-card, .simulation-card {
  margin-top: 20px;
}

.unit {
  margin-left: 5px;
  color: #606266;
}

.simulation-loading {
  display: flex;
  flex-direction: column;
  align-items: center;
  justify-content: center;
  height: 400px;
}

.no-simulation {
  display: flex;
  align-items: center;
  justify-content: center;
  height: 400px;
  color: #909399;
  font-size: 16px;
}

.simulation-summary {
  margin-top: 20px;
  padding-top: 20px;
  border-top: 1px solid #ebeef5;
}

.simulation-summary h3 {
  margin-top: 0;
  margin-bottom: 15px;
  font-size: 16px;
  color: #303133;
}

.summary-item {
  background-color: #f5f7fa;
  padding: 15px;
  border-radius: 4px;
}

.summary-label {
  font-size: 14px;
  color: #606266;
  margin-bottom: 5px;
}

.summary-value {
  font-size: 18px;
  font-weight: bold;
  color: #303133;
}

.summary-item.warning .summary-value {
  color: #e6a23c;
}
</style>
\end{lstlisting}

\subsection{6.2
最佳实践与模式}\label{ux6700ux4f73ux5b9eux8df5ux4e0eux6a21ux5f0f}

在智慧水利平台的Vue.js开发中，一些最佳实践和设计模式可以帮助提高代码质量和开发效率。

\subsubsection{6.2.1
组件设计与复用}\label{ux7ec4ux4ef6ux8bbeux8ba1ux4e0eux590dux7528}

智慧水利平台中的组件设计应遵循单一职责原则，创建可复用的组件体系。

\textbf{组件层次结构}：

\begin{enumerate}
\def\labelenumi{\arabic{enumi}.}
\tightlist
\item
  \textbf{基础UI组件}：封装基本的UI元素，如按钮、输入框、选择器等
\item
  \textbf{业务组件}：封装特定业务逻辑的组件，如水位图表、流量监测卡片等
\item
  \textbf{页面组件}：组合多个业务组件，实现完整的页面功能
\end{enumerate}

\textbf{组件通信模式}：

\begin{lstlisting}
<!-- 父组件传递配置到子组件 -->
<template>
  <div class="watershed-analysis">
    <h2>{{ watershedName }}流域分析</h2>
    
    <water-quality-chart 
      :station-ids="selectedStations"
      :date-range="dateRange"
      :indicators="selectedIndicators"
      :chart-type="chartType"
      @data-loaded="handleDataLoaded"
      @chart-click="handleChartClick">
    </water-quality-chart>
    
    <div class="chart-controls">
      <el-select v-model="chartType" placeholder="选择图表类型">
        <el-option label="线图" value="line"></el-option>
        <el-option label="柱状图" value="bar"></el-option>
        <el-option label="雷达图" value="radar"></el-option>
      </el-select>
      
      <!-- 其他控制选项 -->
    </div>
  </div>
</template>

<script>
import WaterQualityChart from '@/components/charts/WaterQualityChart.vue'

export default {
  components: {
    WaterQualityChart
  },
  data() {
    return {
      watershedName: '长江',
      selectedStations: ['ST001', 'ST002', 'ST003'],
      dateRange: {
        start: new Date(Date.now() - 30 * 24 * 60 * 60 * 1000),
        end: new Date()
      },
      selectedIndicators: ['ph', 'do', 'cod', 'bod'],
      chartType: 'line',
      chartData: null
    }
  },
  methods: {
    handleDataLoaded(data) {
      this.chartData = data;
      console.log('图表数据加载完成:', data);
    },
    handleChartClick(params) {
      console.log('图表点击事件:', params);
      // 处理图表点击，如显示详情等
    }
  }
}
</script>
\end{lstlisting}

\subsubsection{6.2.2
状态管理策略}\label{ux72b6ux6001ux7ba1ux7406ux7b56ux7565}

大型智慧水利平台应用需要合理的状态管理策略，确保数据流转清晰。

\textbf{模块化状态设计}：

\begin{lstlisting}
// store/index.js
import Vue from 'vue'
import Vuex from 'vuex'
import createPersistedState from 'vuex-persistedstate'

// 水利平台各模块状态
import auth from './modules/auth'
import monitoring from './modules/monitoring'
import reservoir from './modules/reservoir'
import forecast from './modules/forecast'
import alert from './modules/alert'
import settings from './modules/settings'

Vue.use(Vuex)

export default new Vuex.Store({
  modules: {
    auth,         // 认证状态
    monitoring,   // 监测数据
    reservoir,    // 水库信息
    forecast,     // 预报信息
    alert,        // 预警信息
    settings      // 用户设置
  },
  plugins: [
    // 持久化用户认证和设置状态
    createPersistedState({
      key: 'water-platform-state',
      paths: ['auth.token', 'auth.user', 'settings']
    })
  ]
})
\end{lstlisting}

\textbf{水文站点数据模块示例}：

\begin{lstlisting}
// store/modules/monitoring.js
import { StationAPI } from '@/api/station'

const state = {
  stations: [],
  realtimeData: {},
  historicalData: {},
  dataLoading: false,
  selectedStationId: null,
  dateRange: {
    start: null,
    end: null
  }
}

const getters = {
  // 获取特定状态的站点
  stationsByStatus: (state) => (status) => {
    return state.stations.filter(station => {
      const data = state.realtimeData[station.id]
      if (!data) return false
      
      if (status === 'warning') {
        return data.waterLevel > station.warningLevel && data.waterLevel <= station.dangerLevel
      } else if (status === 'danger') {
        return data.waterLevel > station.dangerLevel
      } else if (status === 'normal') {
        return data.waterLevel <= station.warningLevel
      }
      
      return false
    })
  },
  
  // 获取选中站点的数据
  selectedStationData: (state) => {
    if (!state.selectedStationId) return null
    
    const station = state.stations.find(s => s.id === state.selectedStationId)
    if (!station) return null
    
    return {
      station,
      realtimeData: state.realtimeData[state.selectedStationId] || null,
      historicalData: state.historicalData[state.selectedStationId] || []
    }
  }
}

const mutations = {
  SET_STATIONS(state, stations) {
    state.stations = stations
  },
  
  UPDATE_REALTIME_DATA(state, { stationId, data }) {
    Vue.set(state.realtimeData, stationId, {
      ...data,
      updateTime: new Date()
    })
  },
  
  SET_HISTORICAL_DATA(state, { stationId, data }) {
    Vue.set(state.historicalData, stationId, data)
  },
  
  SET_SELECTED_STATION(state, stationId) {
    state.selectedStationId = stationId
  },
  
  SET_DATE_RANGE(state, { start, end }) {
    state.dateRange.start = start
    state.dateRange.end = end
  },
  
  SET_LOADING(state, isLoading) {
    state.dataLoading = isLoading
  }
}

const actions = {
  // 获取所有站点基本信息
  async fetchStations({ commit }) {
    commit('SET_LOADING', true)
    
    try {
      const stations = await StationAPI.getStations()
      commit('SET_STATIONS', stations)
    } catch (error) {
      console.error('Failed to fetch stations:', error)
      throw error
    } finally {
      commit('SET_LOADING', false)
    }
  },
  
  // 获取实时数据
  async fetchRealtimeData({ commit, state }, stationIds = null) {
    const ids = stationIds || state.stations.map(s => s.id)
    
    try {
      const dataMap = await StationAPI.getRealtimeData(ids)
      
      Object.entries(dataMap).forEach(([stationId, data]) => {
        commit('UPDATE_REALTIME_DATA', { stationId, data })
      })
    } catch (error) {
      console.error('Failed to fetch realtime data:', error)
      throw error
    }
  },
  
  // 获取历史数据
  async fetchHistoricalData({ commit, state }, { stationId, startTime, endTime }) {
    commit('SET_LOADING', true)
    
    try {
      const data = await StationAPI.getHistoricalData(stationId, startTime, endTime)
      commit('SET_HISTORICAL_DATA', { stationId, data })
    } catch (error) {
      console.error(`Failed to fetch historical data for station ${stationId}:`, error)
      throw error
    } finally {
      commit('SET_LOADING', false)
    }
  },
  
  // 选择站点并加载数据
  async selectStation({ commit, dispatch }, stationId) {
    commit('SET_SELECTED_STATION', stationId)
    
    if (stationId) {
      // 获取最新实时数据
      await dispatch('fetchRealtimeData', [stationId])
      
      // 获取最近7天历史数据
      const endTime = new Date()
      const startTime = new Date(endTime.getTime() - 7 * 24 * 60 * 60 * 1000)
      
      dispatch('fetchHistoricalData', {
        stationId,
        startTime,
        endTime
      })
    }
  }
}

export default {
  namespaced: true,
  state,
  getters,
  mutations,
  actions
}
\end{lstlisting}

\subsubsection{6.2.3
API与服务封装}\label{apiux4e0eux670dux52a1ux5c01ux88c5}

智慧水利平台通常需要与多个后端服务交互，良好的API封装可以提高代码可维护性。

\begin{lstlisting}
// api/request.js - 封装HTTP请求
import axios from 'axios'
import store from '@/store'
import router from '@/router'
import { Message } from 'element-ui'

// 创建axios实例
const service = axios.create({
  baseURL: process.env.VUE_APP_API_BASE_URL,
  timeout: 30000
})

// 请求拦截器
service.interceptors.request.use(
  config => {
    // 添加认证token
    const token = store.state.auth.token
    if (token) {
      config.headers['Authorization'] = `Bearer ${token}`
    }
    
    return config
  },
  error => {
    console.error('Request error:', error)
    return Promise.reject(error)
  }
)

// 响应拦截器
service.interceptors.response.use(
  response => {
    const res = response.data
    
    // 假设API返回格式为 { code: 0, data: {}, message: '' }
    if (res.code !== 0) {
      const errorMessage = res.message || '请求失败'
      Message.error(errorMessage)
      
      // 处理特定错误码
      if (res.code === 401) {
        // 登录过期，重定向到登录页
        store.dispatch('auth/logout')
        router.push('/login')
      }
      
      return Promise.reject(new Error(errorMessage))
    }
    
    return res.data
  },
  error => {
    console.error('Response error:', error)
    
    let errorMessage = '网络错误，请稍后重试'
    if (error.response) {
      const { status, data } = error.response
      
      switch (status) {
        case 401:
          errorMessage = '未授权，请重新登录'
          store.dispatch('auth/logout')
          router.push('/login')
          break
        case 403:
          errorMessage = '拒绝访问'
          break
        case 404:
          errorMessage = '请求的资源不存在'
          break
        case 500:
          errorMessage = '服务器错误'
          break
        default:
          errorMessage = data.message || errorMessage
      }
    }
    
    Message.error(errorMessage)
    return Promise.reject(error)
  }
)

export default service

// api/station.js - 水文站点API
import request from './request'

export const StationAPI = {
  // 获取所有站点信息
  getStations() {
    return request({
      url: '/stations',
      method: 'get'
    })
  },
  
  // 获取站点详情
  getStationById(stationId) {
    return request({
      url: `/stations/${stationId}`,
      method: 'get'
    })
  },
  
  // 获取实时数据
  getRealtimeData(stationIds) {
    return request({
      url: '/stations/realtime',
      method: 'get',
      params: {
        ids: stationIds instanceof Array ? stationIds.join(',') : stationIds
      }
    })
  },
  
  // 获取历史数据
  getHistoricalData(stationId, startTime, endTime) {
    return request({
      url: `/stations/${stationId}/history`,
      method: 'get',
      params: {
        startTime: startTime instanceof Date ? startTime.toISOString() : startTime,
        endTime: endTime instanceof Date ? endTime.toISOString() : endTime
      }
    })
  },
  
  // 添加新站点
  addStation(stationData) {
    return request({
      url: '/stations',
      method: 'post',
      data: stationData
    })
  },
  
  // 更新站点信息
  updateStation(stationId, stationData) {
    return request({
      url: `/stations/${stationId}`,
      method: 'put',
      data: stationData
    })
  }
}

// api/reservoir.js - 水库API
import request from './request'

export const ReservoirAPI = {
  // 获取所有水库信息
  getReservoirs() {
    return request({
      url: '/reservoirs',
      method: 'get'
    })
  },
  
  // 获取水库详情
  getReservoirById(reservoirId) {
    return request({
      url: `/reservoirs/${reservoirId}`,
      method: 'get'
    })
  },
  
  // 获取实时数据
  getRealtimeData(reservoirId) {
    return request({
      url: `/reservoirs/${reservoirId}/realtime`,
      method: 'get'
    })
  },
  
  // 运行调度模拟
  simulateOperation(reservoirId, parameters) {
    return request({
      url: `/reservoirs/${reservoirId}/simulate`,
      method: 'post',
      data: parameters
    })
  },
  
  // 提交调度方案
  submitOperation(reservoirId, operation) {
    return request({
      url: `/reservoirs/${reservoirId}/operations`,
      method: 'post',
      data: operation
    })
  }
}
\end{lstlisting}

\subsection{6.3
性能优化与实用技巧}\label{ux6027ux80fdux4f18ux5316ux4e0eux5b9eux7528ux6280ux5de7}

智慧水利平台通常需要处理大量的地理信息和实时数据，性能优化至关重要。

\subsubsection{6.3.1
大数据量表格优化}\label{ux5927ux6570ux636eux91cfux8868ux683cux4f18ux5316}

智慧水利平台经常需要展示大量的监测数据，可以通过虚拟滚动技术优化表格性能。

\begin{lstlisting}
<!-- VirtualTable.vue -->
<template>
  <div class="virtual-table-container">
    <el-table
      v-loading="loading"
      :data="visibleData"
      :height="height"
      @scroll="handleScroll"
      v-bind="$attrs"
      v-on="$listeners">
      <slot></slot>
    </el-table>
    
    <div class="virtual-placeholder" :style="{ height: `${placeholderHeight}px` }"></div>
  </div>
</template>

<script>
export default {
  name: 'VirtualTable',
  props: {
    data: {
      type: Array,
      required: true
    },
    loading: {
      type: Boolean,
      default: false
    },
    itemHeight: {
      type: Number,
      default: 40
    },
    visibleItemCount: {
      type: Number,
      default: 20
    },
    height: {
      type: Number,
      default: 400
    },
    bufferSize: {
      type: Number,
      default: 5
    }
  },
  data() {
    return {
      startIndex: 0,
      scrollTop: 0
    }
  },
  computed: {
    visibleData() {
      const start = Math.max(0, this.startIndex - this.bufferSize)
      const end = Math.min(this.data.length, this.startIndex + this.visibleItemCount + this.bufferSize)
      
      return this.data.slice(start, end).map((item, index) => ({
        ...item,
        _virtualIndex: start + index
      }))
    },
    totalHeight() {
      return this.data.length * this.itemHeight
    },
    placeholderHeight() {
      return this.totalHeight - this.visibleData.length * this.itemHeight
    }
  },
  methods: {
    handleScroll(e) {
      this.scrollTop = e.target.scrollTop
      this.updateVisibleData()
    },
    updateVisibleData() {
      this.startIndex = Math.floor(this.scrollTop / this.itemHeight)
    }
  }
}
</script>

<style scoped>
.virtual-table-container {
  position: relative;
}

.virtual-placeholder {
  width: 1px;
  visibility: hidden;
}
</style>
\end{lstlisting}

使用示例：

\begin{lstlisting}
<template>
  <div class="data-list">
    <h2>水文站点历史数据</h2>
    
    <div class="filter-bar">
      <el-date-picker 
        v-model="dateRange" 
        type="daterange" 
        range-separator="至" 
        start-placeholder="开始日期" 
        end-placeholder="结束日期"
        @change="handleDateChange">
      </el-date-picker>
      
      <el-select v-model="selectedStation" placeholder="选择站点" @change="handleStationChange">
        <el-option 
          v-for="station in stations" 
          :key="station.id" 
          :label="station.name" 
          :value="station.id">
        </el-option>
      </el-select>
    </div>
    
    <virtual-table 
      :data="historyData" 
      :loading="loading" 
      :height="600"
      :item-height="40"
      :visible-item-count="15"
      border
      stripe>
      <el-table-column type="index" width="50" label="#"></el-table-column>
      <el-table-column prop="time" label="时间" width="180"></el-table-column>
      <el-table-column prop="waterLevel" label="水位(米)" width="120"></el-table-column>
      <el-table-column prop="flowRate" label="流量(m³/s)" width="120"></el-table-column>
      <el-table-column prop="precipitation" label="降水量(mm)" width="120"></el-table-column>
      <el-table-column prop="temperature" label="温度(°C)" width="120"></el-table-column>
      <el-table-column prop="humidity" label="湿度(%)" width="120"></el-table-column>
      <el-table-column label="水位状态" width="120">
        <template slot-scope="scope">
          <el-tag :type="getWaterLevelStatusType(scope.row)">
            {{ getWaterLevelStatusText(scope.row) }}
          </el-tag>
        </template>
      </el-table-column>
    </virtual-table>
  </div>
</template>

<script>
import VirtualTable from '@/components/common/VirtualTable.vue'
import { StationAPI } from '@/api/station'

export default {
  components: {
    VirtualTable
  },
  data() {
    return {
      stations: [],
      selectedStation: '',
      dateRange: [
        new Date(new Date().getTime() - 30 * 24 * 60 * 60 * 1000),
        new Date()
      ],
      historyData: [],
      loading: false
    }
  },
  async created() {
    try {
      this.stations = await StationAPI.getStations()
      if (this.stations.length > 0) {
        this.selectedStation = this.stations[0].id
        this.loadHistoryData()
      }
    } catch (error) {
      console.error('Failed to load stations:', error)
    }
  },
  methods: {
    async loadHistoryData() {
      if (!this.selectedStation) return
      
      this.loading = true
      try {
        const [startDate, endDate] = this.dateRange
        this.historyData = await StationAPI.getHistoricalData(
          this.selectedStation,
          startDate,
          endDate
        )
      } catch (error) {
        console.error('Failed to load history data:', error)
        this.$message.error('加载历史数据失败')
      } finally {
        this.loading = false
      }
    },
    
    handleStationChange() {
      this.loadHistoryData()
    },
    
    handleDateChange() {
      this.loadHistoryData()
    },
    
    getWaterLevelStatusType(row) {
      const station = this.stations.find(s => s.id === this.selectedStation)
      if (!station) return ''
      
      if (row.waterLevel > station.dangerLevel) {
        return 'danger'
      } else if (row.waterLevel > station.warningLevel) {
        return 'warning'
      } else {
        return 'success'
      }
    },
    
    getWaterLevelStatusText(row) {
      const station = this.stations.find(s => s.id === this.selectedStation)
      if (!station) return ''
      
      if (row.waterLevel > station.dangerLevel) {
        return '超警戒'
      } else if (row.waterLevel > station.warningLevel) {
        return '警戒'
      } else {
        return '正常'
      }
    }
  }
}
</script>
\end{lstlisting}

\subsubsection{6.3.2
数据可视化优化}\label{ux6570ux636eux53efux89c6ux5316ux4f18ux5316}

智慧水利平台中的图表数据可视化需要注意性能优化：

\begin{enumerate}
\def\labelenumi{\arabic{enumi}.}
\tightlist
\item
  \textbf{按需加载ECharts组件}：减小包体积
\item
  \textbf{数据下采样}：对于长时间序列数据，可以进行下采样
\item
  \textbf{分段加载}：超大数据集采用分段请求和加载策略
\end{enumerate}

\begin{lstlisting}
// utils/chart-optimizations.js
/**
 * 时间序列数据下采样
 * @param {Array} data 原始数据数组
 * @param {Number} maxPoints 最大点数
 * @returns {Array} 下采样后的数据
 */
export function downsampleTimeSeries(data, maxPoints = 1000) {
  if (!data || data.length <= maxPoints) {
    return data;
  }
  
  const result = [];
  const step = Math.ceil(data.length / maxPoints);
  
  for (let i = 0; i < data.length; i += step) {
    // 在每个步长区间内找最大、最小和平均值进行保留
    const segment = data.slice(i, Math.min(i + step, data.length));
    
    // 如果是单个数据点，直接添加
    if (segment.length === 1) {
      result.push(segment[0]);
      continue;
    }
    
    // 获取区间内的时间范围
    const startTime = segment[0].time;
    const endTime = segment[segment.length - 1].time;
    
    // 找出区间内的最大值和最小值
    let maxValue = -Infinity;
    let minValue = Infinity;
    let maxIndex = 0;
    let minIndex = 0;
    let sum = 0;
    
    segment.forEach((point, index) => {
      const value = point.value;
      sum += value;
      
      if (value > maxValue) {
        maxValue = value;
        maxIndex = index;
      }
      
      if (value < minValue) {
        minValue = value;
        minIndex = index;
      }
    });
    
    // 添加开始点、最小值点、最大值点和结束点
    // 按照时间顺序添加，避免线条交叉
    const pointsToAdd = new Set([0, minIndex, maxIndex, segment.length - 1]);
    
    // 转换为数组并排序
    const sortedIndices = [...pointsToAdd].sort((a, b) => a - b);
    
    // 添加到结果中
    sortedIndices.forEach(index => {
      result.push(segment[index]);
    });
  }
  
  return result;
}

/**
 * 分段加载大数据集
 * @param {Function} loadFunction 数据加载函数
 * @param {String} startTime 开始时间
 * @param {String} endTime 结束时间
 * @param {Number} segments 分段数量
 * @returns {Promise<Array>} 完整数据集
 */
export async function loadLargeDatasetInSegments(loadFunction, startTime, endTime, segments = 4) {
  const startDate = new Date(startTime).getTime();
  const endDate = new Date(endTime).getTime();
  const timeRange = endDate - startDate;
  const segmentSize = timeRange / segments;
  
  const promises = [];
  
  for (let i = 0; i < segments; i++) {
    const segmentStart = new Date(startDate + i * segmentSize);
    const segmentEnd = new Date(Math.min(startDate + (i + 1) * segmentSize, endDate));
    
    promises.push(loadFunction(segmentStart.toISOString(), segmentEnd.toISOString()));
  }
  
  try {
    const results = await Promise.all(promises);
    // 合并所有分段数据
    return results.flat().sort((a, b) => new Date(a.time) - new Date(b.time));
  } catch (error) {
    console.error('Failed to load data segments:', error);
    throw error;
  }
}
\end{lstlisting}

\subsubsection{6.3.3
WebSocket实时数据更新}\label{websocketux5b9eux65f6ux6570ux636eux66f4ux65b0}

智慧水利平台中的实时数据监测可以使用WebSocket进行更高效的更新：

\begin{lstlisting}
// services/websocket.js
import store from '@/store'

export default class WebSocketService {
  constructor(url) {
    this.url = url
    this.socket = null
    this.reconnectInterval = 5000 // 重连间隔(ms)
    this.reconnectAttempts = 0
    this.maxReconnectAttempts = 5
    this.heartbeatTimer = null
    this.heartbeatInterval = 30000 // 心跳间隔(ms)
    this.listeners = {}
    this.connected = false
  }
  
  connect() {
    if (this.socket) {
      this.disconnect()
    }
    
    try {
      this.socket = new WebSocket(this.url)
      
      this.socket.onopen = () => {
        console.log('WebSocket连接已建立')
        this.connected = true
        this.reconnectAttempts = 0
        this.startHeartbeat()
        
        // 发送认证信息
        const token = store.state.auth.token
        if (token) {
          this.send({
            type: 'auth',
            token
          })
        }
        
        this.emit('connect')
      }
      
      this.socket.onmessage = (event) => {
        try {
          const data = JSON.parse(event.data)
          
          // 处理不同类型的消息
          switch (data.type) {
            case 'stationUpdate':
              // 更新站点实时数据
              store.commit('monitoring/UPDATE_REALTIME_DATA', {
                stationId: data.stationId,
                data: data.data
              })
              break
              
            case 'alert':
              // 处理预警信息
              store.dispatch('alert/addAlert', data.alert)
              break
              
            case 'heartbeat':
              // 处理心跳响应
              break
              
            default:
              // 触发自定义事件
              this.emit(data.type, data)
          }
        } catch (error) {
          console.error('Failed to parse WebSocket message:', error)
        }
      }
      
      this.socket.onerror = (error) => {
        console.error('WebSocket连接错误:', error)
        this.emit('error', error)
      }
      
      this.socket.onclose = (event) => {
        this.connected = false
        this.stopHeartbeat()
        console.log(`WebSocket连接已关闭，代码: ${event.code}，原因: ${event.reason}`)
        this.emit('disconnect', event)
        
        // 尝试重连
        if (this.reconnectAttempts < this.maxReconnectAttempts) {
          this.reconnectAttempts++
          console.log(`尝试重连... (${this.reconnectAttempts}/${this.maxReconnectAttempts})`)
          
          setTimeout(() => {
            this.connect()
          }, this.reconnectInterval)
        }
      }
    } catch (error) {
      console.error('Failed to create WebSocket connection:', error)
    }
  }
  
  disconnect() {
    if (this.socket) {
      this.stopHeartbeat()
      this.socket.close()
      this.socket = null
      this.connected = false
    }
  }
  
  send(data) {
    if (!this.socket || this.socket.readyState !== WebSocket.OPEN) {
      console.error('WebSocket is not connected')
      return false
    }
    
    try {
      this.socket.send(JSON.stringify(data))
      return true
    } catch (error) {
      console.error('Failed to send WebSocket message:', error)
      return false
    }
  }
  
  startHeartbeat() {
    this.stopHeartbeat()
    
    this.heartbeatTimer = setInterval(() => {
      if (this.socket && this.socket.readyState === WebSocket.OPEN) {
        this.send({ type: 'heartbeat' })
      }
    }, this.heartbeatInterval)
  }
  
  stopHeartbeat() {
    if (this.heartbeatTimer) {
      clearInterval(this.heartbeatTimer)
      this.heartbeatTimer = null
    }
  }
  
  on(event, callback) {
    if (!this.listeners[event]) {
      this.listeners[event] = []
    }
    
    this.listeners[event].push(callback)
    return this
  }
  
  off(event, callback) {
    if (!this.listeners[event]) return this
    
    if (callback) {
      this.listeners[event] = this.listeners[event].filter(cb => cb !== callback)
    } else {
      delete this.listeners[event]
    }
    
    return this
  }
  
  emit(event, ...args) {
    if (!this.listeners[event]) return
    
    this.listeners[event].forEach(callback => {
      try {
        callback(...args)
      } catch (error) {
        console.error(`Error in WebSocket event listener (${event}):`, error)
      }
    })
  }
  
  // 订阅站点实时数据更新
  subscribeToStations(stationIds) {
    return this.send({
      type: 'subscribe',
      entity: 'stations',
      ids: stationIds
    })
  }
  
  // 订阅水库实时数据更新
  subscribeToReservoirs(reservoirIds) {
    return this.send({
      type: 'subscribe',
      entity: 'reservoirs',
      ids: reservoirIds
    })
  }
  
  // 订阅预警信息
  subscribeToAlerts(areaCode = null) {
    return this.send({
      type: 'subscribe',
      entity: 'alerts',
      areaCode
    })
  }
}

// 创建和配置WebSocket服务实例
const wsService = new WebSocketService(process.env.VUE_APP_WEBSOCKET_URL)

// 在Vue组件中使用
export function useWebSocket() {
  return {
    connect: () => wsService.connect(),
    disconnect: () => wsService.disconnect(),
    subscribeToStations: (ids) => wsService.subscribeToStations(ids),
    subscribeToReservoirs: (ids) => wsService.subscribeToReservoirs(ids),
    subscribeToAlerts: (areaCode) => wsService.subscribeToAlerts(areaCode),
    on: (event, callback) => wsService.on(event, callback),
    off: (event, callback) => wsService.off(event, callback)
  }
}
\end{lstlisting}

在组件中使用WebSocket订阅实时数据：

```vue

\begin{lstlisting}
<h1>智慧水利实时监测</h1>

<el-row :gutter="20">
  <el-col :span="24">
    <el-card>
      <div slot="header">
        <span>重点水文站点实时水位</span>
        <el-tag type="success" size="small" v-if="connected">实时连接中</el-tag>
        <el-tag type="danger" size="small" v-else>未连接</el-tag>
      </div>
      
      <div class="station-grid">
        <station-card 
          v-for="station in keyStations" 
          :key="station.id"
          :station="station"
          :realtime-data="stationData[station.id]">
        </station-card>
      </div>
    </el-card>
  </el-col>
</el-row>

<!-- 更多监测内容 -->
\end{lstlisting}
